\scanpage{69}
% archive sheet number: 68
% @ L1–2
A) Recollements donnant lieu à une application naturelle
\[
\mathscr X\longrightarrow X,
\]
% @ L3–8
ou ce qui revient au $=$, où l'isom $\sigma:\alpha\longrightarrow\beta$ \inserted{(\deleted{\unsure{avec|moy.}} sans composer avec l'antipodisme)} est obtenu par transport parallèle -- ou aussi, combinatoirement, cas où l'isom de transition $\omega_D\longrightarrow\omega_{D'}$ est opposée au transport parallèle. Il y a deux possibilités (avec tjrs $\beta=b$)

% @ L9–13
1) Cas $\alpha=a'$ -- on \unsure{vrille|recolle} \deleted{l'arc de \unsure{sp}} \deleted{une rive $b$ de} l'arc de spéc $b$ sur $\Delta_D$ avec l'arc \deleted{$a'$ qui est} de gén. $a'$ sur $\Delta_{D'}$ qui est sur la rive \uline{faisant face} à celle de $b$, i.e opposée à celle obtenue par transport parallèle. \deleted{C'est} Alors $\mathscr X\to X$ est une immersion. \deleted{On repère}

% @ L14–19
\deleted{\unsure{C|On repère}} Pour définir la structure polygonale sur l'ens des transitions : à partir d'une p.v. donnée $D'$, on répère \deleted{\unsure{l'une des|d'une face}} \inserted{une op} de générisation par la rive qui est \deleted{\unsure{dans la direc|dans la direction}} celle vers laquelle on va déplacer $D'$ pour génériser : le trait bleu qui marque la rive ``tire'' sur $D'$ pour la déplacer vers lui.

% @ L20–21 to R6
2) Cas $\alpha=a$, -- on recolle $b$ avec l'arc $a$ qui s'en déduit par transport parallèle. Maintenant $\mathscr X\to X$ est une immersion sur chaque face $F^*$, mais c'est un rebroussement le long d'une frontière commune à deux faces.

% @ R7–17
On répère une op. de générisation par \deleted{la rive \unsure{li}} l'arc de gén. qui est opposé à la rive vers laquelle on se \unsure{pousse|dirige} -- le trait bleu marquant cette rive va venir ``pousser'' sur $D'$ pour la génériser, \deleted{et} la \unsure{détachant|détacher|direction} d'un sommet envisagé
